% GENERATED FILE. Edit canonical lessons, then run npm run generate:books.
% canonical-source-hash: fnv1a64:77e4a98504735017
% canonical-lessons: IT-C417-banchina, IT-C417-pensilina, IT-C417-capolinea, IT-C417-pullman, IT-C417-transitare

\chapter{Alla fermata}
\label{ch:it-alla-fermata}
\hlchaptermodality{417}

\begin{chapteropening}
\textbf{By the end of this chapter:} \emph{I can wait for a bus or train: the platform, the shelter, the terminus, the coach and trains passing through.}

Complete the last lesson of chapter 417: I can wait for a bus or train: the platform, the shelter, the terminus, the coach and trains passing through.
\end{chapteropening}

\section[la banchina]{la banchina — the platform; the quay}

\label{lesson:IT-C417-banchina}



\begin{quote}
Before the new one: say the Italian for the parking meter, then the Italian for to stop briefly.
\end{quote}

\subsection*{You'll want to know: la banchina}
\textbf{la banchina} — \textquotedblleft{}the platform; the quay\textquotedblright{}.

\begin{quote}
Aspettiamo sulla banchina del binario tre.
\end{quote}

\begin{quote}
\emph{We wait on the platform of track three.}
\end{quote}

\emph{Il binario} is the track and its number; \emph{la banchina} is where you stand. At a port it is the quay.

\begin{cousinweb}
\emph{banchina} is a little \emph{banco} 'bench, counter', which was borrowed from a Germanic word for a bench (EtymDB).
\end{cousinweb}

\begin{grammarlens}[title={What you've built}]
One more word for telling, explaining and describing at B1.
\end{grammarlens}

\subsection*{Your turn}
\begin{itemize}
\raggedright
  \item \emph{Say it:} \emph{la banchina}
  \item \emph{Say it:} \emph{la banchina}, once more
  \item \emph{Say it:} a sentence of your own with \emph{la banchina}
  \item \emph{Recall:} say the Italian for the parking meter, then the Italian for to stop briefly, then say \emph{la banchina} again
\end{itemize}

\begin{tcolorbox}[breakable,colback=teal!4,colframe=teal!35!black,title={Before you move on}]
What does la banchina mean? (\textquotedblleft{}The platform; the quay\textquotedblright{}.) Say it once more.
\end{tcolorbox}
\section[la pensilina]{la pensilina — the bus shelter, the platform canopy}

\label{lesson:IT-C417-pensilina}



\begin{quote}
Before the new one: say the Italian for to stop briefly, then the Italian for the platform.
\end{quote}

\subsection*{You'll want to know: la pensilina}
\textbf{la pensilina} — \textquotedblleft{}the bus shelter, the platform canopy\textquotedblright{}.

\begin{quote}
Ci siamo riparati sotto la pensilina.
\end{quote}

\begin{quote}
\emph{We took shelter under the bus shelter.}
\end{quote}

The roof over a bus stop or a platform.

\begin{cousinweb}
\emph{pensilina} is built on \emph{pensile} 'hanging', a learned borrowing from Latin \emph{pēnsilis} 'hanging' (EtymDB), from \emph{pendēre} 'to hang': a roof that hangs out over you.
\end{cousinweb}

\begin{grammarlens}[title={What you've built}]
One more word for telling, explaining and describing at B1.
\end{grammarlens}

\subsection*{Your turn}
\begin{itemize}
\raggedright
  \item \emph{Say it:} \emph{la pensilina}
  \item \emph{Say it:} \emph{la pensilina}, once more
  \item \emph{Say it:} a sentence of your own with \emph{la pensilina}
  \item \emph{Recall:} say the Italian for to stop briefly, then the Italian for the platform, then say \emph{la pensilina} again
\end{itemize}

\begin{tcolorbox}[breakable,colback=teal!4,colframe=teal!35!black,title={Before you move on}]
What does la pensilina mean? (\textquotedblleft{}The bus shelter, the platform canopy\textquotedblright{}.) Say it once more.
\end{tcolorbox}
\section[il capolinea]{il capolinea — the terminus, the end of the line}

\label{lesson:IT-C417-capolinea}



\begin{quote}
Before the new one: say the Italian for the platform, then the Italian for the bus shelter.
\end{quote}

\subsection*{You'll want to know: il capolinea}
\textbf{il capolinea} — \textquotedblleft{}the terminus, the end of the line\textquotedblright{}.

\begin{quote}
Scendo al capolinea.
\end{quote}

\begin{quote}
\emph{I get off at the terminus.}
\end{quote}

Plural \emph{i capolinea}. Figuratively, \emph{siamo al capolinea} — this is the end.

\begin{cousinweb}
\emph{capolinea} is built from \emph{capo} 'head, end', inherited from Latin \emph{caput} 'head', and \emph{linea} 'line', from Latin \emph{līnea} 'linen thread' (EtymDB).
\end{cousinweb}

\begin{grammarlens}[title={What you've built}]
One more word for telling, explaining and describing at B1.
\end{grammarlens}

\subsection*{Your turn}
\begin{itemize}
\raggedright
  \item \emph{Say it:} \emph{il capolinea}
  \item \emph{Say it:} \emph{il capolinea}, once more
  \item \emph{Say it:} a sentence of your own with \emph{il capolinea}
  \item \emph{Recall:} say the Italian for the platform, then the Italian for the bus shelter, then say \emph{il capolinea} again
\end{itemize}

\begin{tcolorbox}[breakable,colback=teal!4,colframe=teal!35!black,title={Before you move on}]
What does il capolinea mean? (\textquotedblleft{}The terminus, the end of the line\textquotedblright{}.) Say it once more.
\end{tcolorbox}
\section[il pullman]{il pullman — the long-distance bus, the coach}

\label{lesson:IT-C417-pullman}



\begin{quote}
Before the new one: say the Italian for the bus shelter, then the Italian for the terminus.
\end{quote}

\subsection*{You'll want to know: il pullman}
\textbf{il pullman} — \textquotedblleft{}the long-distance bus, the coach\textquotedblright{}.

\begin{quote}
Il pullman per il mare parte alle otto.
\end{quote}

\begin{quote}
\emph{The coach to the seaside leaves at eight.}
\end{quote}

For trips between towns; in town you take \emph{l'autobus}.

\begin{cousinweb}
\emph{pullman} was borrowed from English \emph{pullman}, first the name of a kind of comfortable railway carriage; in Italian it moved to the road.
\end{cousinweb}

\begin{grammarlens}[title={What you've built}]
One more word for telling, explaining and describing at B1.
\end{grammarlens}

\subsection*{Your turn}
\begin{itemize}
\raggedright
  \item \emph{Say it:} \emph{il pullman}
  \item \emph{Say it:} \emph{il pullman}, once more
  \item \emph{Say it:} a sentence of your own with \emph{il pullman}
  \item \emph{Recall:} say the Italian for the bus shelter, then the Italian for the terminus, then say \emph{il pullman} again
\end{itemize}

\begin{tcolorbox}[breakable,colback=teal!4,colframe=teal!35!black,title={Before you move on}]
What does il pullman mean? (\textquotedblleft{}The long-distance bus, the coach\textquotedblright{}.) Say it once more.
\end{tcolorbox}
\section[transitare]{transitare — to pass through, to transit}

\label{lesson:IT-C417-transitare}



\begin{quote}
Before the new one: say the Italian for the terminus, then the Italian for the long-distance bus.
\end{quote}

\subsection*{You'll want to know: transitare}
\textbf{transitare} — \textquotedblleft{}to pass through, to transit\textquotedblright{}.

\begin{quote}
Il treno non ferma qui, transita soltanto.
\end{quote}

\begin{quote}
\emph{The train does not stop here, it only passes through.}
\end{quote}

On station boards: \emph{treno in transito} — a train passing through.

\begin{cousinweb}
\emph{transitare} is built on \emph{transito}, a learned borrowing from Latin \emph{trānsitus} 'a passing over', from \emph{trānsīre} 'to go across', from \emph{trāns} and \emph{īre} 'to go' (Lewis \& Short).
\end{cousinweb}

\begin{grammarlens}[title={What you've built}]
One more word for telling, explaining and describing at B1.
\end{grammarlens}

\subsection*{Your turn}
\begin{itemize}
\raggedright
  \item \emph{Say it:} \emph{transitare}
  \item \emph{Say it:} \emph{transitare}, once more
  \item \emph{Say it:} a sentence of your own with \emph{transitare}
  \item \emph{Recall:} say the Italian for the terminus, then the Italian for the long-distance bus, then say \emph{transitare} again
\end{itemize}

\begin{tcolorbox}[breakable,colback=teal!4,colframe=teal!35!black,title={Before you move on}]
What does transitare mean? (\textquotedblleft{}To pass through, to transit\textquotedblright{}.) Say it once more.
\end{tcolorbox}
